\documentclass[preview]{standalone}

\usepackage{amsmath}
\usepackage{amssymb}
\usepackage{stellar}
\usepackage{definitions}

\begin{document}

\id{psicologia-categorie-concetti}
\genpage

\section{Categorie e concetti}

\begin{snippet}{categorie-introduzione}
    L'esperienza del mondo è un flusso continuo e disomogeneo di sensazioni,
    movimenti, pensieri, emozioni e azioni. La mente non può controllare ogni
    elemento singolarmente; per questo suddivide l'esperienza in unità di
    informazione e raggruppa elementi associabili per somiglianza, contiguità o
    altri criteri.
\end{snippet}

\begin{snippetdefinition}{categoria-mentale-definition}{Categoria mentale}
    Una \emph{categoria mentale} è una classe di entità relativamente omogenee
    al proprio interno ed eterogenee rispetto alle entità di altre classi.
\end{snippetdefinition}

\begin{snippetdefinition}{concetto-definition}{Concetto}
    Un \emph{concetto} è la conoscenza che abbiamo di una categoria di oggetti o
    eventi.
\end{snippetdefinition}

\begin{snippetdefinition}{intensione-definition}{Intensione}
    L'\emph{intensione} di un concetto è l'insieme degli attributi necessari
    affinché un oggetto o evento faccia parte di quel concetto.
\end{snippetdefinition}

\begin{snippetdefinition}{estensione-definition}{Estensione}
    L'\emph{estensione} di un concetto comprende tutti gli oggetti o eventi che
    fanno parte del concetto stesso.
\end{snippetdefinition}

\begin{snippet}{approccio-classico-concetti}
    Secondo l'approccio classico, un concetto è caratterizzato da un insieme di
    attributi necessari e sufficienti. I concetti sono organizzati
    gerarchicamente: gli attributi di un concetto subordinato includono quelli
    del concetto sopraordinato.
\end{snippet}

\includesnpt[width=56\%,src=/snippet/static-psychology/hierarchical-concept-network.jpg]{centered-img}

\begin{snippet}{limite-rete-gerarchica}
    L'idea che ogni esemplare rappresenti una categoria nello stesso modo è
    problematica. Per esempio, lo struzzo appartiene alla categoria degli
    uccelli, ma non possiede la proprietà di volare, tipica di molti altri
    uccelli.
\end{snippet}

\begin{snippetdefinition}{tassonomia-definition}{Tassonomia}
    Una \emph{tassonomia} è un ordinamento sistematico. Nel caso delle
    categorie, indica una classificazione gerarchica delle categorie stesse.
\end{snippetdefinition}

\begin{snippet}{livelli-tassonomici-rosch}
    Eleanor Rosch distingue tre livelli categoriali:
    \begin{enumerate}
        \item \emph{livello sovraordinato}, come \emph{arredamento};
        \item \emph{livello di base}, come \emph{sedia};
        \item \emph{livello subordinato}, come \emph{sedia da cucina}.
    \end{enumerate}
    Le categorie di base sono particolarmente rilevanti perché implicano
    programmi motori unitari, somiglianze morfologiche e parole usate
    frequentemente nella comunicazione.
\end{snippet}

\includesnpt[width=60\%,src=/snippet/static-psychology/category-taxonomy-levels.jpg]{centered-img}

\begin{snippetdefinition}{prototipo-definition}{Prototipo}
    Il \emph{prototipo} è il miglior esemplare di una categoria, cioè quello che
    la rappresenta meglio perché possiede maggiore salienza.
\end{snippetdefinition}

\begin{snippet}{teoria-prototipo}
    Secondo Rosch, i prototipi sono gli elementi centrali attorno ai quali si
    organizza una categoria. L'appartenenza categoriale è graduale: più un
    elemento è simile al prototipo, più forte appare la sua appartenenza alla
    categoria.
\end{snippet}

\begin{snippetdefinition}{rappresentativita-definition}{Rappresentatività}
    La \emph{rappresentatività} indica il possesso del maggior numero di
    proprietà tipiche di una categoria.
\end{snippetdefinition}

\begin{snippetdefinition}{appartenenza-categoriale-definition}{Appartenenza categoriale}
    L'\emph{appartenenza categoriale} indica il fatto che un elemento faccia
    effettivamente parte di una categoria.
\end{snippetdefinition}

\begin{snippetexample}{categoria-uccello-example}{Categoria uccello}
    Per appartenere alla categoria \emph{uccello}, un animale deve possedere
    proprietà essenziali come essere oviparo e avere il becco. Un animale può
    somigliare a un uccello per alcuni aspetti, ma restare fuori dalla categoria
    se non possiede tali proprietà.
\end{snippetexample}

\begin{snippetdefinition}{modello-esteso-prototipo-definition}{Modello esteso del prototipo}
    Il \emph{modello esteso del prototipo} considera il prototipo non come un
    esemplare concreto, ma come una configurazione astratta di proprietà salienti
    che distinguono una categoria da un'altra.
\end{snippetdefinition}

\begin{snippetdefinition}{proprieta-essenziali-definition}{Proprietà essenziali}
    Le \emph{proprietà essenziali} sono proprietà che definiscono l'appartenenza
    categoriale in negativo, perché permettono di escludere dalla categoria gli
    elementi che non le possiedono.
\end{snippetdefinition}

\begin{snippetdefinition}{proprieta-tipiche-definition}{Proprietà tipiche}
    Le \emph{proprietà tipiche} sono proprietà frequenti e salienti di una
    categoria, ma soggette a eccezioni e non indispensabili per l'appartenenza.
\end{snippetdefinition}

\begin{snippetexample}{proprieta-uccelli-example}{Proprietà essenziali e tipiche degli uccelli}
    Per la categoria \emph{uccello}, proprietà essenziali sono essere oviparo e
    avere il becco. Sono invece proprietà tipiche la capacità di volare, avere
    piume nel senso comune del termine e avere ali: alcune specie, come struzzi
    o pinguini, non possiedono tutte queste proprietà tipiche.
\end{snippetexample}

\end{document}
